\section{Sheaves of algebraic structures}
\label{section-sheaves-structures}

\noindent
Let us clarify the definition of sheaves of certain types of structures.
First, let us reformulate the sheaf condition. Namely, suppose that
$\mathcal{F}$ is a presheaf of sets on the topological space $X$.
The sheaf condition can be reformulated as follows. Let
$U = \bigcup_{i\in I} U_i$ be an open covering. Consider the
diagram
$$
\xymatrix{
\mathcal{F}(U) \ar[r]
&
\prod\nolimits_{i\in I}
\mathcal{F}(U_i)
\ar@<1ex>[r] \ar@<-1ex>[r]
&
\prod\nolimits_{(i_0, i_1) \in I \times I}
\mathcal{F}(U_{i_0} \cap U_{i_1})
}
$$
Here the left map is defined by the rule
$s \mapsto \prod_{i \in I} s|_{U_i}$. The two maps
on the right are the maps
$$
\prod\nolimits_i s_i
\mapsto
\prod\nolimits_{(i_0, i_1)} s_{i_0}|_{U_{i_0} \cap U_{i_1}}
\text{ resp. }
\prod\nolimits_i s_i
\mapsto
\prod\nolimits_{(i_0, i_1)} s_{i_1}|_{U_{i_0} \cap U_{i_1}}.
$$
The sheaf condition exactly says that the left arrow
is the equalizer of the right two. This generalizes
immediately to the case of presheaves with values in a
category as long as the category has products.

\begin{definition}
\label{definition-sheaf-values-in-category}
Let $X$ be a topological space. Let $\mathcal{C}$ be
a category with products. A presheaf $\mathcal{F}$ with
values in $\mathcal{C}$ on $X$ is a {\it sheaf}
if for every open covering the diagram
$$
\xymatrix{
\mathcal{F}(U) \ar[r]
&
\prod\nolimits_{i\in I}
\mathcal{F}(U_i)
\ar@<1ex>[r] \ar@<-1ex>[r]
&
\prod\nolimits_{(i_0, i_1) \in I \times I}
\mathcal{F}(U_{i_0} \cap U_{i_1})
}
$$
is an equalizer diagram in the category $\mathcal{C}$.
\end{definition}

\noindent
Suppose that $\mathcal{C}$ is a category and that
$F : \mathcal{C} \to \textit{Sets}$ is a faithful functor.
A good example to keep in mind is the case where $\mathcal{C}$
is the category of abelian groups and $F$ is the forgetful functor.
Consider a presheaf $\mathcal{F}$ with values in $\mathcal{C}$ on $X$.
We would like to reformulate the condition above in terms
of the underlying presheaf of sets
(Definition \ref{definition-underlying-presheaf-sets}).
Note that the underlying
presheaf of sets is a sheaf of sets if and only if all the
diagrams
$$
\xymatrix{
F(\mathcal{F}(U)) \ar[r]
&
\prod\nolimits_{i\in I}
F(\mathcal{F}(U_i))
\ar@<1ex>[r] \ar@<-1ex>[r]
&
\prod\nolimits_{(i_0, i_1) \in I \times I}
F(\mathcal{F}(U_{i_0} \cap U_{i_1}))
}
$$
of sets -- after applying the forgetful functor $F$ -- are
equalizer diagrams! Thus we would like $\mathcal{C}$ to have
products and equalizers and we would like $F$ to commute with
them. This is equivalent to the condition that $\mathcal{C}$
has limits and that $F$ commutes with them, see
Categories, Lemma \ref{categories-lemma-limits-products-equalizers}.
But this is not yet good enough
(see Example \ref{example-sheaves-topological-spaces});
we also need $F$ to {\it reflect isomorphisms}.
This property means that given a morphism
$f : A \to A'$ in $\mathcal{C}$, then $f$ is
an isomorphism if (and only if) $F(f)$ is a bijection.

\begin{lemma}
\label{lemma-sheaves-structure}
Suppose the category $\mathcal{C}$ and
the functor $F : \mathcal{C} \to \textit{Sets}$
have the following properties:
\begin{enumerate}
\item $F$ is faithful,
\item $\mathcal{C}$ has limits and $F$ commutes with them, and
\item the functor $F$ reflects isomorphisms.
\end{enumerate}
Let $X$ be a topological space. Let $\mathcal{F}$
be a presheaf with values in $\mathcal{C}$.
Then $\mathcal{F}$ is a sheaf if and only if the
underlying presheaf of sets is a sheaf.
\end{lemma}

\begin{proof}
Assume that $\mathcal{F}$ is a sheaf. Then
$\mathcal{F}(U)$ is the equalizer of the diagram
above and by assumption we see $F(\mathcal{F}(U))$
is the equalizer of the corresponding diagram
of sets. Hence $F(\mathcal{F})$ is a sheaf of sets.

\medskip\noindent
Assume that $F(\mathcal{F})$ is a sheaf.
Let $E \in \Ob(\mathcal{C})$ be the
equalizer of the two parallel arrows in
Definition \ref{definition-sheaf-values-in-category}.
We get a canonical morphism $\mathcal{F}(U) \to E$,
simply because $\mathcal{F}$ is a presheaf.
By assumption, the induced map $F(\mathcal{F}(U)) \to F(E)$
is an isomorphism, because $F(E)$ is the equalizer
of the corresponding diagram of sets. Hence we
see $\mathcal{F}(U) \to E$ is an isomorphism
by condition (3) of the lemma.
\end{proof}

\noindent
The lemma in particular applies to
{\it sheaves of groups, rings, algebras over a fixed ring, modules
over a fixed ring, vector spaces over a fixed field, } etc.
In other words, these are presheaves of groups, rings,
modules over a fixed ring, vector spaces over a fixed field, etc
such that the underlying presheaf of sets is a sheaf.

\begin{example}
\label{example-C0-sheaf-rings}
Let $X$ be a topological space. For each open $U \subset X$ consider
the $\mathbf{R}$-algebra
$\mathcal{C}^{0}(U) = \{ f : U \to \mathbf{R} \mid f\text{ is continuous}\}$.
There are obvious restriction mappings that turn this into a
presheaf of $\mathbf{R}$-algebras over $X$.
By Example \ref{example-basic-continuous-maps} it is a sheaf of sets.
Hence by the Lemma \ref{lemma-sheaves-structure} it is a sheaf of
$\mathbf{R}$-algebras over $X$.
\end{example}

\begin{example}
\label{example-sheaves-topological-spaces}
Consider the category of topological spaces $\textit{Top}$.
There is a natural faithful functor $\textit{Top} \to \textit{Sets}$
which commutes with products and equalizers. But it does
not reflect isomorphisms. And, in fact it turns out that
the analogue of Lemma \ref{lemma-sheaves-structure} is wrong.
Namely, suppose $X = \mathbf{N}$ with the discrete
topology. Let $A_i$, for $i \in \mathbf{N}$ be a discrete
topological space. For any subset $U \subset \mathbf{N}$
define $\mathcal{F}(U) = \prod_{i\in U} A_i$ with the
discrete topology. Then this is a presheaf of topological
spaces whose underlying presheaf of sets is a sheaf, see
Example \ref{example-sheaf-product-pointwise}.
However, if each $A_i$ has at least two elements, then
this is not a sheaf of topological spaces
according to Definition \ref{definition-sheaf-values-in-category}.
The reader may check that putting the {\it product topology} on
each $\mathcal{F}(U) = \prod_{i\in U} A_i$ does lead to a sheaf
of topological spaces over $X$.
\end{example}


